\clearpage\SourceLeaf{346}

\begin{center}
\textbf{[DIE KEGELSCHNITTE.]}
\end{center}

\begin{center}
[I.]

\textbf{LEHRSATZ.}
\end{center}

Sind $t$, $t'$ die complexen Werthe zweier conjugirter Halbaxen einer
Ellipse und $e$ der complexe Werth der Excentricität (Abstand eines
Brennpunkts), so ist
\[
tt + t't' = ee.
\]

Die Punkte einer Ellipse sind durch die Formel gegeben:
\[
g e^{ix} + h e^{-ix},
\]
wo $x$ alle Werthe von $x = 0$ bis $x = 2\pi$ durchlaufen muss.

Lehrsatz. Es sind dann zwei conjugirte Halbaxen
\[
g e^{ix} + h e^{-ix} \quad \text{und} \quad g i e^{ix} - h i e^{-ix}.
\]
Also
\[
4gh = ee.
\]
